\section{Institutional Background}
\label{sec:institutional}

\subsection{Health as a Constitutional Right and the Rise of Litigation}
\label{sec:health_right}

Brazil's 1988 Constitution created the Unified Health System (SUS), which maintains a formulary of approved medications subject to cost-effectiveness criteria \citep{castro2019, soares2019}. Courts, however, have adopted a strict literal interpretation of Article~196 (``health is a right of all and a duty of the state''), under which any individual can obtain virtually any medication through litigation---from common analgesics to rare-disease treatments costing over US\$\BPrareDrugCostUSD{} per year. Access is inexpensive (a prescription and a public defender suffice); approximately \BPlitGrantedShare{} of first-instance claims are granted \citep{cnj2019judicializacao}. The result is \BPlitClaimsBR{} health court orders nationally in \BPcnjClaimsYear, a \BPspGrowthSince{} increase in S\~ao Paulo since \BPspGrowthBaseYear{} (Online Appendix, Figure~A.1). Litigated health spending in S\~ao Paulo alone amounts to approximately US\$\BPlitSpendingUSDmm~million annually---nearly \BPlitSpendingShare{} of the state's public health budget.

\subsection{The Sanction Channel: Administrative vs.\ Litigated Purchases}
\label{sec:sanction_channel}

The S\~ao Paulo State Department of Health (SES/SP) manages health policy through \BPpbuCount{} decentralized public buyer units (PBUs). All purchases---ordinary and urgent---are conducted through the \textit{Bolsa Eletr\^onica de Compras} (BEC), S\~ao Paulo's electronic procurement platform, mandatory for common goods since \BPbecMandatorySinceYear. Two competitive procedures are used: sealed-bid tendering (\textit{convite}) and multi-round descending auctions (\textit{preg\~ao}). For ordinary purchases, the internal phase typically spans \BPordPlanLow--\BPordPlanHigh~days; almost all court decisions (\BPinjunctionShare) are enforced as preliminary injunctions with delivery deadlines of \BPurgentPlanLow--\BPurgentPlanHigh~days, compressing the internal phase to roughly one-third of the ordinary timeline.

A critical institutional feature for our analysis is the distinction between \textit{litigated} and \textit{administrative} urgent purchases. Both types face identical planning constraints---compressed timelines, small quantities, diverted budget resources, and the same BEC platform and bidding procedures. They differ sharply in the consequences of procurement failure. For litigated purchases, failure to comply with a court order exposes public officials to substantial fines (sometimes disproportionate to the purchase value), administrative and criminal liability, and seizure or blocking of public funds; these penalties are public information, as the court order must be referenced in the tender notice. Since \BPsampleStartYear, the SES/SP has run a parallel \textit{administrative-request} channel: a mechanism allowing the government to negotiate supply directly with individuals before litigation occurs. Administrative requests undergo evaluation by a scientific committee using evidence-based medicine criteria, and---crucially---failure to complete an administrative purchase carries no penalties for public officials. Between \BPsampleStartYear{} and \BPsampleEndYear, administrative requests generated approximately \BPadminPurchasesTotal{} purchases, or about \BPadminToCourtShare{} of the total from court orders.

This channel structure is the institutional vehicle for our identification. Administrative and litigated urgent purchases share all execution constraints but differ only in the penalty regime, so a within-item comparison between them is the cleanest available measurement of the cost margin attributable to personal sanctions on procurement officials. Figure~\ref{fig:purchase_types} summarizes the three purchase types. We discuss the selection of cases into the administrative channel---and the threats it poses to a strict causal interpretation---in Section~\ref{sec:empirical}.

\begin{table}[ht]
  \centering
  \caption{Purchase Types: Institutional Characteristics}
  \label{fig:purchase_types}
  \small
  \begin{threeparttable}
  \begin{tabular}{lccc}
    \toprule
    & Ordinary & Administrative & Litigated \\
    \midrule
    Source of funds     & Dedicated budget & Diverted budget & Diverted budget \\
    Quantity            & Large            & Small           & Small \\
    Delivery time       & Long (\BPordPlanLow--\BPordPlanHigh{} d) & Short (\BPurgentPlanLow--\BPurgentPlanHigh{} d) & Short (\BPurgentPlanLow--\BPurgentPlanHigh{} d) \\
    Threat of punishment & None            & None            & \textbf{Fines, liability,} \\
                        &                  &                 & \textbf{fund seizure} \\
    \bottomrule
  \end{tabular}
  \begin{tablenotes}
    \small
    \item \textit{Notes:} Administrative and litigated purchases share planning
    constraints (short timelines, small quantities, diverted budgets) and differ
    only in the penalty regime. The ``under the gun'' comparison exploits this
    variation.
  \end{tablenotes}
  \end{threeparttable}
\end{table}
